\documentclass[12pt,reqno]{amsart}
\usepackage{math-review}
% Real Banach spaces; all time norms on a fixed finite interval.
\newcommand{\trans}{\tau}
\newcommand{\cone}{\mathcal M_{+}}
\newcommand{\gauge}[1]{[#1]_{\cone}}
\newcommand{\R}{\mathbb R}

\title{Time compactness: which hypotheses can be transferred?}
\date{}
\begin{document}
\maketitle
\section{Main Results}
The question is which spatial and time estimates can be transferred without changing the unknown or the convergence norm. The full statements below retain their decisive assumptions before any comparison is made.
\begin{theorem}[Simon]\label{thm:simon}
Let $T>0$, $1\le p\le\infty$, and let $X,B,Y$ be Banach spaces with continuous linear injections
\[
X\Subset B\hookrightarrow Y,
\]
where the first injection is compact. Let $F\subset L^p(0,T;X)$ satisfy
\[
\sup_{f\in F}\|f\|_{L^p(0,T;X)}<\infty,
\qquad
\lim_{h\downarrow0}\sup_{f\in F}\|f(\cdot+h)-f\|_{L^p(0,T-h;Y)}=0.
\]
Then $F$ is relatively compact in $L^p(0,T;B)$ for $p<\infty$.
For $p=\infty$, the $B$-valued continuous representatives of the elements of $F$ form a relatively compact set in $C([0,T];B)$.
\end{theorem}
The statement is \cite[Theorem 5, pp.~84--85]{S-simon}. The topology used for time control is $Y$, whereas the convergence sought is in $B$.

\begin{theorem}[Maitre]\label{thm:maitre}
Let $T>0$, $1\le p\le\infty$, and let $E_1,E_2$ be Banach spaces.
Let $K:E_1\to E_2$ map every bounded subset of $E_1$ into a relatively compact subset of $E_2$.
Let $U\subset L^1(0,T;E_1)$ be bounded and suppose that each $v(t)=K(u(t))$, $u\in U$, is strongly measurable.
Write $V=\{K(u):u\in U\}$ and assume $V\subset L^p(0,T;E_2)$, that for some $1<r\le\infty$
\[
\sup_{v\in V}\|v\|_{L^r(0,T;E_2)}<\infty,
\qquad
\lim_{h\downarrow0}\sup_{v\in V}\|v(\cdot+h)-v\|_{L^p(0,T-h;E_2)}=0.
\]
Then $V$ is relatively compact in $L^p(0,T;E_2)$ for $p<\infty$.
For $p=\infty$, the continuous representatives form a relatively compact set in $C([0,T];E_2)$.
No injectivity, linearity or monotonicity of $K$ is required. The conclusion concerns $K(u)$, not $u$.
\end{theorem}
This is \cite[Theorem 2.1, p.~1726]{S-maitre}, with the operator renamed $K$. The auxiliary exponent $r$ need not equal $p$. Its role in the argument below is to make a truncated output small in $L^1$.

\begin{theorem}[Barrett--S\"uli]\label{thm:cone}
Let $T>0$, let $A_0,A_1$ be real Banach spaces with a continuous linear injection $A_0\hookrightarrow A_1$, and let $\mathcal M_+\subset A_0$ be closed under multiplication by nonnegative scalars.
Suppose its gauge $[\cdot]:\mathcal M_+\to[0,\infty)$ satisfies
\[
[z]=0\Longleftrightarrow z=0,\qquad [cz]=c[z]\quad(c\ge0),
\qquad \|z\|_{A_0}\le C_0[z],
\]
and every gauge-bounded sequence in $\mathcal M_+$ has a subsequence converging in $A_0$.
Fix $1\le p,q\le\infty$ with either $q>1$, or $q=1$ and $p<\infty$.
Let $F$ be a family of strongly measurable $A_0$-valued functions $v$ on $(0,T)$ with $v(t)\in\mathcal M_+$ almost everywhere and measurable $t\mapsto[v(t)]$.
Assume the distributional derivatives $v'$ are $A_1$-valued and
\[
\sup_{v\in F}\left(\|[v(\cdot)]\|_{L^p(0,T)}+
\|v'\|_{L^q(0,T;A_1)}\right)<\infty.
\]
Then $F$ is relatively compact in $L^p(0,T;A_0)$.
The assertion includes $p=\infty$ only when $q>1$; it does not assert that limits belong to a nonclosed cone.
\end{theorem}
This is the bounded-family form of \cite[Theorem 1 and Remark 2, pp.~3--4]{S-barrett-suli}. A cone need not be closed under subtraction: the displayed axioms do not make $[u-v]$ meaningful. The exceptional pair $(p,q)=(\infty,1)$ is genuinely excluded.
All time intervals have finite length. Relative compactness means that every sequence has a norm-convergent subsequence. We write $\tau_h f(t)=f(t+h)$ on $(0,T-h)$; all translation limits are uniform over the family. Bochner integration, H\"older and Minkowski inequalities, and the fundamental theorem for $W^{1,1}$ trajectories are assumed. The compactness criterion and the nonlinear constructions needed here are stated or developed below.
\section{Key Mathematical Ingredients}
Two tempting stronger conclusions fail for explicit reasons. The examples fix the unknown and endpoint before the usable compactness interfaces are introduced.
\subsection{The excluded endpoint can fail already on a ray}
On $(0,1)$ let $n\ge3$ and define
\[
v_n(t)=\begin{cases}
0,&t\le\tfrac12,\\
n(t-\tfrac12),&\tfrac12<t<\tfrac12+\tfrac1n,\\
1,&t\ge\tfrac12+\tfrac1n.
\end{cases}
\]
These functions take values in the compactly embedded finite-dimensional cone $[0,\infty)\subset\R$ with gauge $[z]=|z|$, and
\[
\|v_n\|_\infty=1,\qquad \|v_n'\|_1=n\cdot n^{-1}=1.
\]
Their pointwise limit is the discontinuous step $w=\mathbf1_{(1/2,1]}$. No subsequence converges uniformly, since a uniform limit of continuous functions is continuous. Continuous-function differences have the same supremum and essential-supremum norms, so there is no $L^\infty$-convergent subsequence either. Yet for every finite $p$,
\begin{equation}\label{eq:ramp}
\|v_n-w\|_p^p=\int_0^{1/n}(1-ns)^p\,ds=\frac1{n(p+1)}\longrightarrow0.
\end{equation}
Thus the distinction between finite $p$ and $p=\infty$ survives the simplest spatial geometry; compare \cite[Corollary 4, pp.~85--86; Proposition 3, p.~94]{S-simon} and \cite[Remark 2, p.~4]{S-barrett-suli}.
\subsection{An image can converge while its input stays separated}
Take $E_1=E_2=\R$, $K(x)=x^2$, and
\[
u_n(t)=(-1)^{\lfloor 2^nt\rfloor},\qquad 0<t<1.
\]
For $m>n$, the signs differ on exactly half the interval. Consequently
\[
\|u_n-u_m\|_{L^1}=2\cdot\tfrac12=1,
\qquad K(u_n(t))=1.
\]
The inputs have no $L^1$-convergent subsequence, while their images have every $L^r$ bound and zero translations. This tests the exact unknown in \cite[Theorem 2.1 and Remark 2.2(3), p.~1726]{S-maitre}. It also shows why the inverse-gauge substitution in Proposition~\ref{prop:overlap} cannot simply be applied to every compact map.
\subsection{What a derivative estimate actually gives}
For $v\in W^{1,1}(0,T;E)$ choose a time with norm no larger than its time average, up to an arbitrarily small error. Integrating $v'$ from that time gives
\begin{equation}\label{eq:w11}
\|v\|_{L^\infty(E)}\le T^{-1}\|v\|_{L^1(E)}+\|v'\|_{L^1(E)}.
\end{equation}
If these supremum and derivative bounds are $K_0,D$, Fubini and bounded amplitudes give, for $p<\infty$,
\begin{align}
\|\tau_hv-v\|_{L^1(E)}
&\le\int_0^{T-h}\int_t^{t+h}\|v'(s)\|_E\,ds\,dt\le hD,\label{eq:l1shift}\\
\|\tau_hv-v\|_{L^p(E)}^p
&\le(2K_0)^{p-1}\|\tau_hv-v\|_{L^1(E)}
\le(2K_0)^{p-1}hD.\label{eq:pshift}
\end{align}
For $q>1$ a different estimate is available, directly from H\"older on $(t,t+h)$:
\begin{equation}\label{eq:supshift}
\|v(t+h)-v(t)\|_E\le h^{1-1/q}\|v'\|_{L^q(E)}.
\end{equation}
Here $1/q=0$ when $q=\infty$. At $q=1$ the exponent is zero, so this last estimate gives no small supremum modulus. These are the interfaces behind \cite[Lemma 4, pp.~77--78]{S-simon} and \cite[Corollary 2.4, p.~1730]{S-maitre}.
\subsection{The output required from spatial information}
\begin{lemma}[Simon]\label{lem:criterion}
Let $E$ be a Banach space, $T>0$, $1\le p\le\infty$, and $F\subset L^p(0,T;E)$.
Assume that for every $0<a<b<T$ the set
\[
\left\{\int_a^b f(t)\,dt:f\in F\right\}
\]
is relatively compact in $E$, and that
\[
\lim_{h\downarrow0}\sup_{f\in F}\|f(\cdot+h)-f\|_{L^p(0,T-h;E)}=0.
\]
Then $F$ is relatively compact in $L^p(0,T;E)$ if $p<\infty$.
If $p=\infty$, its elements have continuous representatives on $[0,T]$ and these representatives form a relatively compact subset of $C([0,T];E)$ with the supremum norm.
\end{lemma}
We import this precise criterion from \cite[Theorem 1, pp.~71--73]{S-simon}. Its proof uses time averages and a reversed-time argument to cover both endpoints. Thus the task is to make every interior interval integral compact, then supply translations in the desired time exponent. For $p=\infty$, the lemma also supplies continuous representatives; arbitrary compact subsets of $L^\infty$ need not satisfy its translation hypothesis.
\subsection{Comparing two points without subtracting in the cone}
\begin{lemma}[Barrett--S\"uli]\label{lem:ehrling}
Let $A_0\hookrightarrow A_1$ be a continuous linear injection of real Banach spaces.
Let $\mathcal M_+\subset A_0$ be a seminormed nonnegative cone with gauge $[\cdot]$, satisfying $\|z\|_{A_0}\le C[z]$ and compactness of gauge-bounded sequences in $A_0$.
Then for every $\eta>0$ there is $C_\eta<\infty$ such that
\[
\|u-v\|_{A_0}\le\eta([u]+[v])+C_\eta\|u-v\|_{A_1}
\quad\text{for all }u,v\in\mathcal M_+.
\]
Here the cone is closed under nonnegative scalar multiplication and its gauge is positive definite and positively degree-one homogeneous; no addition or subtraction closure is assumed.
\end{lemma}
To see the mechanism in \cite[Lemma 1, pp.~4--5]{S-barrett-suli}, suppose the inequality fails at $\eta_0>0$. Choose $u_n,v_n$ with
\[
\|u_n-v_n\|_{A_0}>\eta_0([u_n]+[v_n])+n\|u_n-v_n\|_{A_1}.
\]
The sum $a_n=[u_n]+[v_n]$ is positive. Divide both points by $a_n$. Their gauges are at most one, so a common subsequence converges in $A_0$. The normalized difference satisfies
\[
\|\widetilde u_n-\widetilde v_n\|_{A_0}>\eta_0,
\qquad
\|\widetilde u_n-\widetilde v_n\|_{A_1}
< n^{-1}\|\widetilde u_n-\widetilde v_n\|_{A_0}\longrightarrow0.
\]
Continuity and injectivity of $A_0\hookrightarrow A_1$ force the two $A_0$ limits to agree, contradicting the first inequality. The proof uses scaling of each cone point, never a gauge of their difference.
\section{Connecting the Ingredients}
A valid transfer must preserve the compact object, the norm of the time estimate and its uniform quantifiers. We now perform those checks rather than infer an equivalence from the names of the theorems.
\subsection{Why compact value sets give compact averages}
Let $Q$ be compact in a Banach space $E$. Given $\delta>0$, choose a finite $\delta$-net $q_1,\ldots,q_J$ for $Q$. Replacing each point in a convex combination by a nearby $q_j$ moves that combination by at most $\delta$. Hence
\[
\operatorname{co}(Q)\subset \operatorname{co}\{q_1,\ldots,q_J\}+\delta B_E.
\]
The finite-dimensional simplex on the right is compact. It follows that $\operatorname{co}(Q)$ is totally bounded, so its closure is compact by completeness of $E$.
For a strongly measurable $w:I\to Q$, partition $I$ into measurable nearest-net cells. The corresponding simple-function averages are convex combinations of the $q_j$ and approximate $|I|^{-1}\int_Iw$ within $\delta$. Therefore
\[
|I|^{-1}\int_I w(t)\,dt\in\overline{\operatorname{co}}(Q).
\]
This elementary construction is used for both nonlinear images and gauge-truncated cone values. The averages may lie outside the original image or cone; only their compactness in the ambient Banach space is needed.
\subsection{The nonlinear route: truncate before averaging}
Use the uniform bounds
\[
C_U=\sup_{u\in U}\|u\|_{L^1(E_1)},\qquad
C_V=\sup_{v\in V}\|v\|_{L^r(E_2)}.
\]
For $M>0$, set
\[
G_M(u)=\{t:\|u(t)\|_{E_1}>M\},\quad
u_M=u\,\mathbf1_{(0,T)\setminus G_M(u)},\quad v_M=K(u_M).
\]
Markov's inequality and H\"older give
\begin{align}
|G_M(u)|&\le C_U/M,\label{eq:badset}\\
\|v-v_M\|_{L^1(E_2)}
&\le (C_U/M)^{1-1/r}C_V+(C_U/M)\|K(0)\|_{E_2}.
\label{eq:tail}
\end{align}
The $K(0)$ term is needed because truncation sends the input to zero, not necessarily the output. The right side tends to zero uniformly in $u$. This is an $L^1$ estimate, which remains useful even when the desired $p$ is larger than $r$.
For fixed $M$, the values of $v_M$ lie in the compact set
\[
Q_M=\overline{K(\{x:\|x\|_{E_1}\le M\})}.
\]
The compact-average construction makes $\{\int_Iv_M\}$ relatively compact. Equation~\eqref{eq:tail} approximates $\int_Iv$ uniformly by these integrals. Choosing $M$ first and then a finite net proves that $\{\int_Iv\}$ is totally bounded in $E_2$. Lemma~\ref{lem:criterion} uses the separately assumed $L^p(E_2)$ translations and concludes Theorem~\ref{thm:maitre}.
Maitre's original proof uses shifted finite Riemann sums \cite[Lemma 2.3, equations (2.3)--(2.24), pp.~1726--1730]{S-maitre}. The compact-convex-hull proof here is an elementary reconstruction of the same truncation interface. Neither argument commutes $K$ with an integral.
\subsection{A proof hypothesis is not automatically a sharp boundary}
In \eqref{eq:tail}, $r>1$ gives the positive exponent $1-1/r$. At $r=1$, this particular estimate loses its small factor. That observation does not prove failure of nonlinear compactness at $r=1$.
Later nonlinear translation refinements, including Chen--Liu's \emph{Two nonlinear compactness theorems in $L^p(0,T;B)$}, Applied Mathematics Letters 25 (2012), 2252--2257, lie outside this account. The omission challenge therefore forbids a strongest-known or necessity claim for the auxiliary $r>1$ hypothesis. The endpoint ramp has a different force: it proves actual failure of the claimed $L^\infty$ conclusion under an $L^1$ derivative bound.
\subsection{The linear route: compactness first in the weaker norm}
Let $M=\sup_{f\in F}\|f\|_{L^p(X)}$. For every interior interval $I$,
\[
\left\|\int_I f(t)\,dt\right\|_X\le |I|^{1-1/p}M.
\]
The integrals are relatively compact in $B$ and hence in $Y$. Lemma~\ref{lem:criterion}, used in $Y$, produces a subsequence Cauchy in $L^p(Y)$ (uniformly in $Y$ at $p=\infty$).
The linear Ehrling estimate of \cite[Lemma 8, pp.~84--85]{S-simon} is
\[
\|z\|_B\le\eta\|z\|_X+C_\eta\|z\|_Y,\qquad z\in X.
\]
It can be checked by normalizing a violating sequence to unit $X$ norm: compactness gives a $B$ limit, while the violation forces its $Y$ norm to zero and its $B$ norm away from zero. Injectivity gives the contradiction.
Minkowski now supplies the decisive upgrade:
\begin{equation}\label{eq:linearupgrade}
\|f_n-f_m\|_{L^p(B)}\le 2\eta M+C_\eta\|f_n-f_m\|_{L^p(Y)}.
\end{equation}
For $M>0$, fix $\eta<\varepsilon/(4M)$, then take $n,m$ large enough for the second term to be below $\varepsilon/2$. If $M=0$ the family vanishes. Completeness gives convergence in $L^p(B)$. The continuous-representative assertion at $p=\infty$ is the corresponding part of Simon's theorem, not a property assumed of arbitrary equivalence classes.
\subsection{The cone route: create weak-space compactness, then upgrade}
Let $G=\sup_{v\in F}\|[v]\|_{L^p}$, $D=\sup_{v\in F}\|v'\|_{L^q(A_1)}$, and $C_1=\|A_0\hookrightarrow A_1\|$. The gauge bound and H\"older give
\[
\|v\|_{L^1(A_1)}\le C_1C_0T^{1-1/p}G.
\]
Equation~\eqref{eq:w11} gives a uniform $A_1$ amplitude
\begin{equation}\label{eq:coneamp}
K_1=C_1C_0T^{-1/p}G+T^{1-1/q}D.
\end{equation}
The $A_1$ representatives are continuous, while cone membership is needed only almost everywhere.
Truncate now by the gauge: put $E_R(v)=\{t:[v(t)]>R\}$ and $v_R=v\mathbf1_{E_R(v)^c}$. Then
\begin{equation}\label{eq:conetail}
|E_R(v)|\le\frac{T^{1-1/p}G}{R},\qquad
\|v-v_R\|_{L^1(A_1)}\le\frac{K_1T^{1-1/p}G}{R}.
\end{equation}
The values of $v_R$ have gauge at most $R$, so their closure is compact in $A_0$ and $A_1$. Compact averages and the uniform tail establish compactness of every interval integral in $A_1$.
If $q=1$ and $p<\infty$, equation~\eqref{eq:pshift} gives the required $L^p(A_1)$ translations. If $q>1$, equation~\eqref{eq:supshift} gives supremum translations. Lemma~\ref{lem:criterion} therefore supplies an $L^p(A_1)$-Cauchy subsequence in all permitted cases.
Applying Lemma~\ref{lem:ehrling} pointwise and then Minkowski yields
\begin{equation}\label{eq:coneupgrade}
\|v_n-v_m\|_{L^p(A_0)}\le2\eta G+C_\eta\|v_n-v_m\|_{L^p(A_1)}.
\end{equation}
Choose $\eta$ first, then $n,m$, exactly as in \eqref{eq:linearupgrade}. Completeness proves Theorem~\ref{thm:cone}, including the $L^\infty$ conclusion when $q>1$. If $G=0$, all trajectories vanish almost everywhere.
For clarity at $p=\infty$, one can also recover continuous $A_0$ representatives. On a full-measure dense set of times with $[v(t)]\le G$, Lemma~\ref{lem:ehrling} bounds time differences by $2\eta G+C_\eta\|v(t)-v(s)\|_{A_1}$. First reduce $\eta$, then $|t-s|$; completeness extends the resulting uniformly continuous $A_0$ trajectory to $[0,T]$. The $A_1$ representative agrees by injectivity.
The last upgrade is the one in \cite[equation (2.22), pp.~12--13]{S-barrett-suli}. The preceding weak-space argument is a local reconstruction: it avoids the unresolved dependence of the selected index on time in the passage following that source's (2.8), p.~6. No assertion that the theorem is false follows from this scoped proof-reading concern.
\subsection{An overlap with an explicit change of unknown}
\begin{proposition}\label{prop:overlap}
Let $E_1,E_2$ be real Banach spaces, $T>0$, and $\alpha\ge1$.
Suppose $K:E_1\to E_2$ is injective, maps bounded sets into relatively compact sets, and satisfies $K(cu)=c^\alpha K(u)$ for $c\ge0$.
Set $\mathcal M_+=K(E_1)$ and $[v]=\|K^{-1}v\|_{E_1}^{\alpha}$.
Let $1\le p,q\le\infty$, with $q>1$ or with $q=1,p<\infty$.
Suppose a family $F$ of strongly measurable $\mathcal M_+$-valued trajectories has strongly measurable inverse lifts $u=K^{-1}v$, measurable gauges, and
\[
\sup_{v\in F}\bigl(\|[v]\|_{L^p(0,T)}+\|v'\|_{L^q(0,T;E_2)}\bigr)<\infty,
\]
where derivatives are $E_2$-valued distributions.
Then the hypotheses of the compact nonlinear-map theorem stated for $K$, $U=K^{-1}F$, and $V=F$ hold, using any fixed $r>1$ and the target exponent $p$.
Consequently $F$ is relatively compact in $L^p(0,T;E_2)$, and in $C([0,T];E_2)$ when $q>1$.
This reduction concerns the special cone $K(E_1)$ with the same output and derivative space $E_2$; it is not an equivalence for arbitrary cones or arbitrary compact maps.
\end{proposition}
This is the special interface discussed in \cite[Section 4, pp.~15--16]{S-barrett-suli}, with inverse-lift measurability made explicit. It can be checked directly.
Injectivity and homogeneity give $K(0)=0$ and
\[
K^{-1}(cv)=c^{1/\alpha}K^{-1}(v),\qquad [cv]=c[v].
\]
Since $K$ has bounded image on the unit ball, scaling gives
\[
\|K(x)\|_{E_2}\le C\|x\|_{E_1}^{\alpha}=C[K(x)].
\]
Gauge-bounded sets come from bounded subsets of $E_1$, so they are relatively compact in $E_2$.
If $G=\sup_{v\in F}\|[v]\|_{L^p}$, then $\alpha p\ge1$ and H\"older gives
\begin{equation}\label{eq:lift}
\|K^{-1}v\|_{L^1(E_1)}
\le T^{1-1/(\alpha p)}G^{1/\alpha};
\end{equation}
for $p=\infty$ the factor is $T$. The assumed strongly measurable lifts thus form a bounded $L^1(E_1)$ family. The gauge and derivative bounds make $v$ uniformly bounded in $W^{1,1}(E_2)$, so \eqref{eq:w11} gives every fixed $L^r(E_2)$ bound, $r>1$. Equations~\eqref{eq:pshift} and \eqref{eq:supshift} supply the exact translations; Theorem~\ref{thm:maitre} applies.
The norm matching matters: the derivative here is in $E_2$, the same space as the image. A derivative only in a larger weak space cannot be renamed an $E_2$ derivative. Nor does a general cone arrive with an injective homogeneous parametrization. This is a verified overlap, not an equivalence of all three theorems.
\paragraph{How to read the limits together.}
The ramp invalidates the $(\infty,1)$ conclusion even with finite-dimensional spatial compactness. The sign functions invalidate recovery of $u$ from compactness of $K(u)$. The homogeneous comparison requires an injective parametrization, a measurable inverse lift and a derivative measured in the image space. These are three different checks; none can replace another. The cone proof, by contrast, can use a weaker derivative space because it explicitly performs the two-point norm upgrade.
\clearpage
\bibliographystyle{plain}
\bibliography{../generated/references}
\end{document}
